\section{Introduction}
\label{sec:introduction}

Generative modeling of discrete structured objects, such as graphs, is a central challenge across many scientific and industrial applications, including drug discovery and materials design.

The key difficulty arises from the combinatorial complexity of the underlying data distribution, where domain-specific constraints and patterns significantly shape high-probability regions.
%
Current state-of-the-art models for unconditional graph generation predominantly rely on discrete diffusion~\citep{vignac2022digress} or continuous-time Markov chains~\citep{campbell2024generative,qinmadeira2024defog}, translating samples from simple reference distributions into complex graph structures. Although effective in generating realistic unconditional samples, these methods struggle when enforcing user-specified properties or constraints, as the guidance typically occurs off-manifold—where property regressors and conditioning functions are inherently less reliable.
%\end{mdframed}

Energy-based models (EBMs) offer a complementary perspective that naturally addresses these limitations~\citep{hinton2002training}. EBMs represent the data distribution implicitly through a single scalar energy function, enabling the direct incorporation of constraints, priors, or property-based objectives at inference time without retraining~\citep{du2019implicit, du2025learning}. Particularly for graph generation, EBMs have the advantage that permutation invariance and graph-structured constraints can be explicitly and elegantly integrated into the energy. Yet, historically, EBMs have struggled with sample quality, primarily due to poor sampling efficiency in discrete, combinatorial domains, as exemplified by existing methods like GraphEBM~\citep{liu2021graphebm}. Consequently, discrete EBMs have predominantly served as scoring models for out-of-distribution detection rather than as competitive generative samplers~\citep{wu2023energy,fuchsgruber2024energy}.


Energy Matching~\citep{balcerak2025energy} bridges these perspectives in continuous spaces by learning a time-independent potential that simultaneously (i) guides samples from a reference distribution toward data, and (ii) defines a Boltzmann-like density near data via contrastive refinement. This regularizes energies off-support while retaining the compositional advantages of energy-based models (EBMs). However, the continuous formulation still relies on explicit, time-indexed sampling schedules (via ordinary or stochastic differential equations; ODEs/SDEs) to trigger exploration and mixing. Extending Energy Matching to discrete graph spaces, we remove this explicit time dependence entirely, instead using local edits within a Markov Chain Monte Carlo (MCMC) sampler. Here, an energy-dependent criterion implicitly switches between exploration (transport) and refinement regimes, naturally matching the combinatorial structure of graphs. Effective graph-based generation thus requires symmetry-aware transport supervision and efficient discrete sampling strategies.

\Gls{jko} provides a variational framework for optimization in spaces of probability measures, where each
step balances an energy against a distance (or discrepancy) rooted in \gls{ot}, such as the Wasserstein
metric~\citep{lanzetti2025firstorder,lanzetti2024variational,terpin2024learning}. This perspective motivates OT-inspired transport supervision under a given geometry.
For graphs, \gls{gw}-type costs offer permutation-respecting
comparisons~\citep{krzakala2024any2graph}. At the same time, computing exact Wasserstein geodesics is
typically impractical at scale. Importantly, stochastic interpolants~\citep{albergoBuildingNormalizingFlows2023} with data-dependent couplings~\citep{albergo2024stochastic} provide a principled way to work with tractable, possibly suboptimal couplings; i.e., minibatch coupling induces practically useful paths that guide sampling toward data-supported regions without requiring exact solutions to optimal transport. This motivates using efficient approximate graph couplings to supervise global guidance toward high-probability regions, combined with local refinement steps to ensure effective exploration and validity.
\begin{figure*}[t]
\centering

% Okabe–Ito palette
\definecolor{otColor}{RGB}{230,159,0}   % orange: OT
\definecolor{mhColor}{RGB}{0,114,178}   % blue:   MH
\tikzset{state/.style={circle,draw,minimum size=11mm,inner sep=0pt}}

\begin{tikzpicture}[>=Stealth, every node/.style={font=\small}, node distance=1.7cm]

% Transport phase nodes (orange shade)
\node[state,fill=orange!20]                    (pi0b) {$\pi_{0}$};
\node[state,fill=orange!20,right=of pi0b]      (pi1b) {$\pi_{1}$};
\node[right=of pi1b]                           (dots1b) {$\dots$};
\node[state,fill=orange!20,right=of dots1b]    (pikb) {$\pi_{k}$};

% Phase-transition divider
\draw[densely dashed,line width=1.2pt] ($(pikb.east)+(0.10,0.95)$) -- ($(pikb.east)+(0.10,-0.95)$);

% Mixing phase nodes (blue shade)
\node[right=1.2cm of pikb]                       (dots2b) {$\dots$};
\node[state,fill=blue!15,right=1.0cm of dots2b]  (pidatab) {$\pi_{\text{data}}$};
\node[right=2mm of pidatab]                      (eqb) {$=$};
\node[state,fill=blue!15,right=2mm of eqb]       (pistb) {$\pi_{\theta}$};

% Transport arrows
\draw[->,very thick,draw=otColor,shorten <=7pt,shorten >=7pt] (pi0b) -- (pi1b);
\draw[->,very thick,draw=otColor,shorten <=7pt,shorten >=7pt] (pi1b) -- (dots1b);
\draw[->,very thick,draw=otColor,shorten <=7pt,shorten >=7pt] (dots1b) -- (pikb);

% MH arrows
\draw[->,very thick,draw=mhColor,shorten <=7pt,shorten >=7pt] (pikb) -- (dots2b);
\draw[->,very thick,draw=mhColor,shorten <=7pt,shorten >=7pt] (dots2b) -- (pidatab);

% Annotation for gradient-informed transport
\node[below=10pt of dots1b,align=center,font=\scriptsize]
      {greedy discrete proposals on samples\\(direction $-\nabla_x V_{\theta}(x)$ $\leftrightarrow$ transport map direction)};


% Annotation for discrete proposals + MH
\node[below=10pt of dots2b,align=center,font=\scriptsize,xshift=36pt]
      {local, geometry-aware, $\nabla_x V_\theta(x)$-informed jumps\\+ MH acceptance};

% Phase annotations
\node[align=center,above=1mm of pi1b]
      {\textbf{Transport Phase}\\[-1pt]\scriptsize samples far from data manifold: $V_{\theta}(x) \gg V_{\text{data}}^{\text{mean}}$};

\node[align=center,above=1mm of pidatab]
      {\textbf{Mixing Phase}\\[-1pt]\scriptsize samples near data manifold: $V_{\theta}(x) \sim V_{\text{data}}^{\text{mean}}$};

% Self-loop at stationary state
\draw[->,thick,draw=mhColor]
  ($(pistb.east)+(2pt,3pt)$) arc[start angle=160,end angle=-160,radius=8pt];

\end{tikzpicture}

\caption{Illustration of GEM sampling, a two-phase MCMC method. Sampling alternates between a transport phase, where gradient-informed, greedy proposals rapidly move samples toward regions of high probability, and a mixing phase employing MH acceptance to ensure correct stationary distribution and efficient mixing between modes. \textbf{Color key:} transport (\textcolor{otColor}{orange}), MH (\textcolor{mhColor}{blue}).}

\label{fig:sampling-regimes}
\end{figure*}


\begin{mdframed}[hidealllines=true,backgroundcolor=blue!5]
\paragraph{Contributions.}
We propose \acrlong{gem}\footnote{Code will be released upon acceptance.}, a state-of-the-art discrete energy-based model for graphs, achieving unconditional generation quality competitive with leading discrete diffusion models. GEM combines:
(i) global transport guidance through permutation-invariant graph couplings under a relaxed OT geometry,
(ii) a learned, time-independent potential explicitly defining a Gibbs distribution over graphs, and
(iii) efficient geometry-aware local edit proposals combined with Metropolis–Hastings sampling.

GEM uniquely retains the compositional and controllable properties of energy-based models, enabling inference-time incorporation of constraints or objectives into the energy. Additionally, its explicit energy structure naturally facilitates computing graph geodesics—minimum-energy edit paths aligned with the learned data geometry—and enables state-of-the-art conditional generation. GEM also achieves state-of-the-art generation of novel graphs by editing existing samples.
\end{mdframed}
